\chapter{Zu wenig und zu viel}

\section{Der Leitfall}

\begin{merksatz}
\boxtitel{Wie schnell sinkt das Laktat?}
Nach Beginn der Sepsistherapie wird das Laktat alle $80$ Minuten gemessen, zehnmal über zwölf Stunden. Gesucht ist eine Kurve, die den Verlauf beschreibt --- um vorherzusagen, wo der Wert zwischen den Messungen liegt.

\begin{center}
{\small
\renewcommand{\arraystretch}{1.36}
\begin{tabular}{@{}P{1.5cm}P{1.4cm}P{0.6cm}P{1.5cm}P{1.4cm}@{}}
\toprule
\textbf{Stunde} & \textbf{Laktat} & & \textbf{Stunde} & \textbf{Laktat}\\
\midrule
$0{,}0$ & $3{,}76$ & & $6{,}7$ & $1{,}60$\\
$1{,}3$ & $2{,}75$ & & $8{,}0$ & $1{,}24$\\
$2{,}7$ & $2{,}47$ & & $9{,}3$ & $1{,}06$\\
$4{,}0$ & $2{,}23$ & & $10{,}7$ & $1{,}43$\\
$5{,}3$ & $2{,}13$ & & $12{,}0$ & $1{,}64$\\
\bottomrule
\end{tabular}}
\end{center}

\schritt{Was man nicht weiß --- aber hier ausnahmsweise weiß}
Die Werte wurden erzeugt aus
\[ L(t) = 1 + 3\,e^{-t/4} \]
plus einem Messfehler mit Streuung $\sigma = 0{,}3$.

\begin{center}
\begin{minipage}{0.88\textwidth}
\kursiv{Das Klinikum kennt die wahre Kurve nicht. \textbf{Wir kennen sie, damit wir prüfen können, wer richtig liegt} --- ein Vorteil, den man bei echten Daten nie hat, und der Grund, warum Lehrbücher des Faches so oft mit erzeugten Daten arbeiten.}
\end{minipage}
\end{center}
\end{merksatz}

\begin{merksatz}
\boxtitel{Drei Versuche}
Man passt Polynome verschiedenen Grades nach der Methode der kleinsten Quadrate an (Band 17, Kapitel~7).

\begin{center}
\begin{tikzpicture}[scale=1.0]
\draw[grau!60, line width=0.6pt] (0,0) -- (7.4,0);
\draw[grau!60, line width=0.6pt] (0,0) -- (0,3.3);
\draw[grau!55, line width=1.0pt, dashed, smooth, samples=60, domain=0:7.2]
  plot ({\x}, {0.8*(1+3*exp(-\x/2.4))-0.6});
\foreach \t/\y in {0.0/3.76,0.8/2.75,1.6/2.47,2.4/2.23,3.2/2.13,4.0/1.60,4.8/1.24,5.6/1.06,6.4/1.43,7.2/1.64}
  \fill[tief] (\t,{0.8*\y-0.6}) circle (2.2pt);
\draw[akzent, line width=1.2pt, smooth] plot coordinates {(0,2.32) (0.8,1.55) (1.6,1.48) (2.4,1.20) (3.2,1.22) (4.0,0.74) (4.8,0.36) (5.6,0.25) (5.9,0.42) (6.4,0.55) (6.8,0.4) (7.2,0.72)};
\draw[gruen!60!black, line width=1.2pt] (0,2.02) -- (7.2,0.22);
\node[font=\scriptsize, color=gruen!60!black, anchor=west] at (7.25,0.22) {Grad 1};
\node[font=\scriptsize, color=akzent, anchor=west] at (7.25,0.85) {Grad 9};
\node[font=\scriptsize, color=grau, anchor=west] at (7.25,1.35) {wahr};
\node[font=\scriptsize\itshape, color=grau] at (3.6,-0.4) {Stunden nach Therapiebeginn};
\end{tikzpicture}
\end{center}

\begin{center}
{\small
\renewcommand{\arraystretch}{1.5}
\begin{tabular}{@{}P{1.6cm}P{2.6cm}P{2.6cm}@{}}
\toprule
\textbf{Grad} & \textbf{Trainingsfehler} & \textbf{bei $t=11$}\\
\midrule
$1$ & $0{,}144$ & $1{,}15$\\
$3$ & $0{,}034$ & $1{,}40$\\
$9$ & $\mathbf{0{,}000}$ & $\mathbf{1{,}83}$\\
\midrule
wahr & --- & $1{,}19$\\
\bottomrule
\end{tabular}}
\end{center}

\begin{center}
\begin{minipage}{0.9\textwidth}
\textbf{Das Polynom neunten Grades läuft durch jeden Messpunkt --- und liegt bei $t=11$ am weitesten daneben.}
\end{minipage}
\end{center}

\begin{center}
\begin{minipage}{0.88\textwidth}
\kursiv{Zehn Punkte legen ein Polynom neunten Grades eindeutig fest; der Trainingsfehler ist also zwangsläufig null. \textbf{Die Kurve folgt dem Messrauschen} --- dem Ausreißer bei $10{,}7$ Stunden ebenso wie dem echten Verlauf --- und schwingt zwischen den Punkten aus. Das ist Aufgabe 1.3 in stetiger Form: Sie hat die Daten gelernt, nicht den Zusammenhang.}
\end{minipage}
\end{center}
\end{merksatz}

\section{Die Zerlegung}

\begin{merksatz}
\boxtitel{Worin ein Vorhersagefehler besteht}
Man denke sich viele Kliniken, die jede ihre eigenen zehn Messungen machen und jede ihr eigenes Polynom anpassen. An einer festen Stelle $x$ gibt jede eine andere Vorhersage $\hat h(x)$ ab.

\schritt{Drei Größen}
\begin{center}
{\small
\renewcommand{\arraystretch}{1.55}
\begin{tabular}{@{}P{2.4cm}P{6.2cm}@{}}
\toprule
\textbf{Rauschen} & $\sigma^{2}$ --- der Messfehler, den kein Verfahren beseitigt.\\
\textbf{Verzerrung} & $\bigl(\E[\hat h(x)] - f(x)\bigr)^{2}$ --- wie weit die \kursiv{mittlere} Vorhersage vom wahren Wert abliegt.\\
\textbf{Varianz} & $\E\bigl[(\hat h(x) - \E[\hat h(x)])^{2}\bigr]$ --- wie stark die Vorhersagen der Kliniken untereinander schwanken.\\
\bottomrule
\end{tabular}}
\end{center}
\end{merksatz}

\begin{merksatz}
\boxtitel{Satz 2.1 --- Verzerrung und Varianz}
Sei $y = f(x) + \varepsilon$ mit $\E[\varepsilon] = 0$ und $\operatorname{Var}[\varepsilon] = \sigma^{2}$, und sei $\varepsilon$ unabhängig von den Trainingsdaten. Dann gilt
\[ \E\bigl[(y - \hat h(x))^{2}\bigr] = \sigma^{2} + \bigl(\E[\hat h(x)] - f(x)\bigr)^{2} + \operatorname{Var}\bigl[\hat h(x)\bigr] . \]

\schritt{Beweis}
Schreibe $\bar h = \E[\hat h(x)]$ und zerlege
\[ y - \hat h = \underbrace{(y-f)}_{\varepsilon} + \underbrace{(f-\bar h)}_{\text{fest}} + \underbrace{(\bar h-\hat h)}_{\text{Schwankung}} . \]

\kursiv{Quadrieren} liefert drei Quadrate und drei gemischte Glieder. Die gemischten Glieder verschwinden im Erwartungswert:

\begin{center}
{\small
\renewcommand{\arraystretch}{1.5}
\begin{tabular}{@{}P{3.0cm}P{5.6cm}@{}}
\toprule
$2\,\varepsilon\,(f-\bar h)$ & $\E[\varepsilon] = 0$, der andere Faktor ist fest.\\
$2\,\varepsilon\,(\bar h-\hat h)$ & $\varepsilon$ unabhängig von den Daten, $\E[\varepsilon]=0$.\\
$2\,(f-\bar h)(\bar h-\hat h)$ & $\E[\bar h-\hat h] = \bar h - \bar h = 0$.\\
\bottomrule
\end{tabular}}
\end{center}

Es bleiben die drei Quadrate:
\[ \E[\varepsilon^{2}] + (f-\bar h)^{2} + \E[(\bar h-\hat h)^{2}] = \sigma^{2} + \text{Verzerrung}^{2} + \text{Varianz} \]
$\blacksquare$

\begin{center}
\begin{minipage}{0.9\textwidth}
\textbf{Der Fehler auf neuen Daten besteht aus drei Teilen --- und nur zwei davon hängen vom Verfahren ab.}
\end{minipage}
\end{center}
\end{merksatz}

\begin{merksatz}
\boxtitel{Der Leitfall zerlegt}
Viertausend simulierte Kliniken, je zehn Messungen, Fehler gemittelt über den ganzen Zeitraum:

\begin{center}
{\small
\renewcommand{\arraystretch}{1.44}
\begin{tabular}{@{}P{1.2cm}P{1.8cm}P{1.8cm}P{1.8cm}P{1.8cm}@{}}
\toprule
\textbf{Grad} & \textbf{Training} & \textbf{Verzerr.$^{2}$} & \textbf{Varianz} & \textbf{Test}\\
\midrule
$0$ & $0{,}886$ & $0{,}606$ & $0{,}009$ & $0{,}705$\\
$1$ & $0{,}184$ & $0{,}081$ & $0{,}017$ & $0{,}187$\\
$\mathbf{2}$ & $0{,}070$ & $0{,}005$ & $0{,}023$ & $\mathbf{0{,}119}$\\
$3$ & $0{,}054$ & $0{,}000$ & $0{,}031$ & $0{,}121$\\
$5$ & $0{,}035$ & $0{,}000$ & $0{,}047$ & $0{,}137$\\
$7$ & $0{,}019$ & $0{,}000$ & $0{,}080$ & $0{,}170$\\
$9$ & $0{,}000$ & $0{,}000$ & $0{,}595$ & $0{,}685$\\
\bottomrule
\end{tabular}}
\end{center}
\kursiv{Die Spalte \enquote{Test} ist $\sigma^{2} + \text{Verzerrung}^{2} + \text{Varianz}$ mit $\sigma^{2} = 0{,}09$.}

\begin{center}
\begin{tikzpicture}[scale=1.0]
\draw[grau!60, line width=0.6pt] (0,0) -- (7.2,0);
\draw[grau!60, line width=0.6pt] (0,0) -- (0,3.7);
\draw[tief, line width=1.2pt] plot coordinates {(0,3.54) (0.8,0.74) (1.6,0.28) (2.4,0.22) (3.2,0.18) (4.0,0.14) (4.8,0.11) (5.6,0.08) (6.4,0.04) (7.2,0.0)};
\draw[akzent, line width=1.4pt] plot coordinates {(0,2.82) (0.8,0.75) (1.6,0.48) (2.4,0.48) (3.2,0.51) (4.0,0.55) (4.8,0.60) (5.6,0.68) (6.4,0.91) (7.2,2.74)};
\fill[akzent] (1.6,0.48) circle (2.6pt);
\node[font=\scriptsize, color=tief, anchor=west] at (4.6,0.35) {Training};
\node[font=\scriptsize, color=akzent, anchor=west] at (5.2,1.35) {Test};
\node[font=\scriptsize, color=grau, anchor=north] at (0,-0.1) {$0$};
\node[font=\scriptsize, color=grau, anchor=north] at (1.6,-0.1) {$2$};
\node[font=\scriptsize, color=grau, anchor=north] at (7.2,-0.1) {$9$};
\node[font=\scriptsize\itshape, color=grau] at (3.6,-0.55) {Grad des Polynoms};
\end{tikzpicture}
\end{center}

\begin{center}
\begin{minipage}{0.9\textwidth}
\textbf{Der Trainingsfehler fällt stetig. Der Testfehler fällt, erreicht bei Grad 2 sein Minimum und steigt wieder.}
\end{minipage}
\end{center}

\schritt{Was man ablesen kann}
\begin{center}
{\small
\renewcommand{\arraystretch}{1.5}
\begin{tabular}{@{}P{2.6cm}P{6.0cm}@{}}
\toprule
\textbf{Grad 0 und 1} & \Begriff{Unteranpassung}: Das Modell ist zu starr, um den Verlauf abzubilden. Die Verzerrung überwiegt.\\
\textbf{Grad 2 und 3} & Beide Anteile klein --- der beste Kompromiss.\\
\textbf{Grad 7 bis 9} & \Begriff{Überanpassung}: Das Modell folgt dem Rauschen. Die Varianz überwiegt.\\
\bottomrule
\end{tabular}}
\end{center}

\begin{center}
\begin{minipage}{0.88\textwidth}
\kursiv{\textbf{Grad 9 ist fast so schlecht wie Grad 0} --- $0{,}685$ gegen $0{,}705$. Wer alles lernt, ist kaum besser dran als wer nichts lernt, und aus genau demselben Grund wie in Aufgabe 1.3. Bemerkenswert ist auch, dass die wahre Kurve gar kein Polynom ist; kein Grad erreicht die Verzerrung null exakt, aber ab Grad 3 ist sie vernachlässigbar.}
\end{minipage}
\end{center}
\end{merksatz}

\section{Wie optimistisch der Trainingsfehler ist}

\begin{merksatz}
\boxtitel{Satz 2.2 --- der Optimismus}
Werden $p$ Parameter nach kleinsten Quadraten an $n$ Messpunkte angepasst und an \kursiv{denselben} Punkten mit neuen Messungen geprüft, so gilt
\[ \E[\text{Testfehler}] - \E[\text{Trainingsfehler}] = \frac{2\,\sigma^{2}\,p}{n} . \]

\schritt{Beweis}
Die Anpassung ist eine lineare Abbildung der Messwerte: $\hat{\vek y} = H\vek y$ mit der Projektionsmatrix $H$ auf den Raum der Polynome vom Grad $p-1$ (Band 2, Kapitel~9).

\kursiv{Erstens:} Der Unterschied zwischen einer neuen Messung $y_i'$ und der alten $y_i$ ist nur das Rauschen. Also
\[ \E\bigl[(y_i'-\hat y_i)^{2}\bigr] - \E\bigl[(y_i-\hat y_i)^{2}\bigr] = 2\operatorname{Cov}(\hat y_i,\,y_i) , \]
denn $y_i'$ ist von $\hat y_i$ unabhängig, $y_i$ aber nicht.

\kursiv{Zweitens:} Da $\hat{\vek y} = H\vek y$ und die Messfehler unabhängig mit Varianz $\sigma^{2}$ sind,
\[ \operatorname{Cov}(\hat y_i,\, y_i) = H_{ii}\,\sigma^{2} . \]

\kursiv{Drittens:} Die Spur einer Projektion ist die Dimension ihres Bildes:
\[ \sum_i H_{ii} = \operatorname{Spur} H = p . \]

\kursiv{Summieren und durch $n$ teilen:}
\[ \frac1n\sum_i 2\,H_{ii}\,\sigma^{2} = \frac{2\,\sigma^{2}\,p}{n} \]
$\blacksquare$

\schritt{Die Probe durch Simulation}
\begin{center}
{\small
\renewcommand{\arraystretch}{1.46}
\begin{tabular}{@{}P{1.2cm}P{1.2cm}P{2.4cm}P{2.4cm}@{}}
\toprule
$n$ & $p$ & \textbf{gemessen} & $2\sigma^{2}p/n$\\
\midrule
$10$ & $2$ & $0{,}0366$ & $0{,}0360$\\
$10$ & $3$ & $0{,}0542$ & $0{,}0540$\\
$10$ & $6$ & $0{,}1072$ & $0{,}1080$\\
$10$ & $10$ & $0{,}1807$ & $0{,}1800$\\
$30$ & $6$ & $0{,}0363$ & $0{,}0360$\\
\bottomrule
\end{tabular}}
\end{center}

\begin{center}
\begin{minipage}{0.9\textwidth}
\textbf{Der Trainingsfehler ist um einen Betrag zu optimistisch, der mit der Zahl der Parameter wächst und mit der Zahl der Beobachtungen fällt.}
\end{minipage}
\end{center}

\begin{center}
\begin{minipage}{0.88\textwidth}
\kursiv{Das ist die Grundlage aller Korrekturen, die man dem Trainingsfehler hinzufügt, um den Testfehler zu schätzen --- das Akaike-Kriterium aus Band 17 ist genau diese Rechnung in anderer Einheit. \textbf{Man \enquote{bezahlt} jeden Parameter mit $2\sigma^{2}/n$.} Kapitel~3 benutzt das, und Kapitel~7 verallgemeinert es auf Verfahren, bei denen man die Parameter nicht zählen kann.}
\end{minipage}
\end{center}

\schritt{Ein Vorbehalt}
\begin{center}
\begin{minipage}{0.88\textwidth}
\kursiv{Der Satz gilt an den Messpunkten selbst. \textbf{Zwischen den Messpunkten kann es erheblich schlimmer sein}: Beim Polynom neunten Grades beträgt der Optimismus an den Punkten $0{,}18$, über den ganzen Zeitraum aber $0{,}685$ --- die Kurve schwingt zwischen den Punkten aus, wo sie nie geprüft wurde. Das ist die zweite Warnung des Vorworts in quantitativer Gestalt.}
\end{minipage}
\end{center}
\end{merksatz}

\section{Mehr Daten}

\begin{merksatz}
\boxtitel{Was mit wachsendem $n$ geschieht}
\begin{center}
{\small
\renewcommand{\arraystretch}{1.46}
\begin{tabular}{@{}P{1.4cm}P{1.6cm}P{1.6cm}P{0.4cm}P{1.6cm}P{1.6cm}@{}}
\toprule
& \multicolumn{2}{c}{\textbf{Grad 2}} & & \multicolumn{2}{c}{\textbf{Grad 9}}\\
$n$ & Verzerr.$^{2}$ & Test & & Varianz & Test\\
\midrule
$10$ & $0{,}005$ & $0{,}118$ & & $0{,}589$ & $0{,}679$\\
$20$ & $0{,}005$ & $0{,}108$ & & $0{,}042$ & $0{,}132$\\
$50$ & $0{,}005$ & $0{,}100$ & & $0{,}017$ & $0{,}107$\\
$100$ & $0{,}004$ & $0{,}097$ & & $0{,}009$ & $0{,}099$\\
\bottomrule
\end{tabular}}
\end{center}

\begin{center}
\begin{minipage}{0.9\textwidth}
\textbf{Mehr Daten senken die Varianz, nicht die Verzerrung.}
\end{minipage}
\end{center}

\begin{center}
\begin{minipage}{0.88\textwidth}
\kursiv{Beim Polynom zweiten Grades bleibt eine Verzerrung von $0{,}004$ bestehen, gleich wie viele Messungen man hat --- die wahre Kurve ist eben keine Parabel. Beim Polynom neunten Grades verschwindet die Varianz, und bei hundert Messungen sind beide fast gleichauf. \textbf{Die richtige Komplexität hängt also von der Datenmenge ab}: Wer wenige Daten hat, muss sich bescheiden; wer viele hat, darf mehr wagen.}
\end{minipage}
\end{center}
\end{merksatz}

\section{Der Bruchfall}

\begin{warnung}
\boxtitelr{Was die Zerlegung voraussetzt}
\begin{center}
{\small
\renewcommand{\arraystretch}{1.52}
\begin{tabular}{@{}P{2.4cm}P{6.2cm}@{}}
\toprule
\textbf{Quadratischer Verlust} & Für den Null-Eins-Verlust aus Kapitel~1 gibt es keine so saubere Zerlegung.\\
\textbf{Festes Rauschen} & $\sigma^{2}$ ist überall gleich. Bei hohem Laktat ist die Messung oft ungenauer.\\
\textbf{Innerhalb des Bereichs} & Geprüft wird zwischen $0$ und $12$ Stunden. Außerhalb gilt nichts.\\
\bottomrule
\end{tabular}}
\end{center}

\schritt{Zur dritten Zeile}
Man frage das Polynom neunten Grades nach dem Wert bei $t = 14$ Stunden.

\begin{center}
\begin{minipage}{0.88\textwidth}
\kursiv{Es liefert eine Zahl, und die Zahl ist absurd --- bei der obigen Messreihe fällt die Kurve jenseits von zwölf Stunden steil ab oder steigt steil an, je nach Vorzeichen des höchsten Koeffizienten. \textbf{Auch das Polynom zweiten Grades wird außerhalb des Bereichs falsch}: Eine Parabel steigt irgendwann wieder, das Laktat nicht. Kein Satz dieses Kapitels sagt etwas über Vorhersagen außerhalb des Bereichs, in dem gemessen wurde.}
\end{minipage}
\end{center}

\schritt{Und ein Wort zum Leitfall}
\begin{center}
\begin{minipage}{0.88\textwidth}
\kursiv{Die zehn Messungen stammen von \kursiv{einem} Patienten. Wer aus ihnen eine Kurve für \kursiv{alle} Patienten ableitet, begeht einen Fehler, den keine Zerlegung erfasst: \textbf{Die Streuung zwischen Patienten ist eine andere Größe als die zwischen Messungen}, und sie ist meist die größere. Band 17, Kapitel~12 behandelt die Modelle, die beides trennen.}
\end{minipage}
\end{center}
\end{warnung}

\section{Übungsaufgaben}

\begin{aufgabe}
\boxtitel{Aufgabe 2.1 --- Die Zerlegung für den Mittelwert}
Man schätzt einen Laktatwert $\mu$ durch den Mittelwert $\bar y$ von $n$ Messungen mit Streuung $\sigma$. Zerlegen Sie den Fehler.

\loesung
\schritt{Schritt 1 --- Die Verzerrung}
\[ \E[\bar y] = \frac1n\sum_{i=1}^{n}\E[y_i] = \frac1n\cdot n\,\mu = \mu \]
\[ \text{Verzerrung} = \E[\bar y] - \mu = 0 \]

\schritt{Schritt 2 --- Die Varianz}
Bei unabhängigen Messungen addieren sich die Varianzen (Band 4, Kapitel~8):
\[ \operatorname{Var}[\bar y] = \frac{1}{n^{2}}\sum_{i=1}^{n}\operatorname{Var}[y_i] = \frac{1}{n^{2}}\cdot n\,\sigma^{2} = \frac{\sigma^{2}}{n} \]

\schritt{Schritt 3 --- Der Gesamtfehler für eine neue Messung}
\[ \E[(y-\bar y)^{2}] = \sigma^{2} + 0 + \frac{\sigma^{2}}{n} = \sigma^{2}\left(1+\frac1n\right) \]

\schritt{Schritt 4 --- Die Probe am Leitfall}
Grad 0 ist genau der Mittelwert. Mit $\sigma^{2} = 0{,}09$ und $n = 10$:
\[ \operatorname{Var} = \frac{0{,}09}{10} = 0{,}009 \]
Die Simulation ergab $0{,}009$. \checkmark

\begin{center}
\begin{minipage}{0.88\textwidth}
\kursiv{\textbf{Der Mittelwert ist unverzerrt --- und trotzdem das schlechteste Verfahren der Tabelle}, weil er den Verlauf gar nicht abbildet. Unverzerrtheit ist eine Eigenschaft bezüglich der Größe, die man schätzt; hier schätzt er eine Konstante, während die Wirklichkeit eine Kurve ist. Die Verzerrung von $0{,}606$ ist keine Verzerrung des Schätzers, sondern des Modells.}
\end{minipage}
\end{center}
\end{aufgabe}

\begin{aufgabe}
\boxtitel{Aufgabe 2.2 --- Absichtlich verzerren}
Ein Laktatwert $\mu = 1{,}2$ wird aus einer einzigen Messung mit Streuung $\sigma = 1$ geschätzt. Man verwendet statt $y$ den geschrumpften Schätzer $c\,y$ mit $0 < c \le 1$. Welches $c$ ist das beste?

\loesung
\schritt{Schritt 1 --- Verzerrung und Varianz}
\[ \E[c\,y] = c\,\mu \quad\Longrightarrow\quad \text{Verzerrung} = (c-1)\,\mu \]
\[ \operatorname{Var}[c\,y] = c^{2}\sigma^{2} \]

\schritt{Schritt 2 --- Der mittlere quadratische Fehler}
\[ M(c) = (c-1)^{2}\mu^{2} + c^{2}\sigma^{2} \]

\schritt{Schritt 3 --- Das Minimum}
\[ M'(c) = 2(c-1)\mu^{2} + 2c\,\sigma^{2} = 0 \]
\[ c\,(\mu^{2}+\sigma^{2}) = \mu^{2} \quad\Longrightarrow\quad c^{*} = \frac{\mu^{2}}{\mu^{2}+\sigma^{2}} \]

\kursiv{Im Beispiel:} $c^{*} = \frac{1{,}44}{1{,}44+1} = \frac{1{,}44}{2{,}44} = 0{,}590$.

\schritt{Schritt 4 --- Die Zahlen}
\begin{center}
{\small
\renewcommand{\arraystretch}{1.48}
\begin{tabular}{@{}P{1.6cm}P{2.0cm}P{1.8cm}P{1.8cm}@{}}
\toprule
$c$ & \textbf{Verzerr.$^{2}$} & \textbf{Varianz} & \textbf{Fehler}\\
\midrule
$1{,}000$ & $0{,}000$ & $1{,}000$ & $1{,}000$\\
$0{,}800$ & $0{,}058$ & $0{,}640$ & $0{,}698$\\
$\mathbf{0{,}590}$ & $0{,}242$ & $0{,}348$ & $\mathbf{0{,}590}$\\
$0{,}400$ & $0{,}518$ & $0{,}160$ & $0{,}678$\\
\bottomrule
\end{tabular}}
\end{center}

\begin{center}
\begin{minipage}{0.9\textwidth}
\textbf{Ein absichtlich verzerrter Schätzer hat einen um vierzig Prozent kleineren Fehler als der unverzerrte.}
\end{minipage}
\end{center}

\schritt{Schritt 5 --- Der Haken}
\begin{center}
\begin{minipage}{0.88\textwidth}
\kursiv{Das beste $c$ hängt von $\mu$ ab --- also von genau der Größe, die man schätzen will. \textbf{Man kann es nicht direkt benutzen.} Aber man sieht, in welche Richtung es geht: Bei großem Rauschen im Verhältnis zum Signal soll man stärker schrumpfen. Kapitel~4 macht daraus ein Verfahren, das mit einer aus den Daten geschätzten Stärke arbeitet --- die Regularisierung.}
\end{minipage}
\end{center}

\schritt{Schritt 6 --- Die Einordnung}
\begin{center}
\begin{minipage}{0.88\textwidth}
\kursiv{In Band 17 galt Unverzerrtheit als Tugend eines Schätzers. \textbf{Hier zeigt sich, dass sie zu teuer sein kann}: Wer um jeden Preis unverzerrt schätzt, zahlt mit Varianz. Das Lernen gibt diese Tugend bewusst auf --- und das ist einer der wenigen Punkte, an denen es der klassischen Statistik tatsächlich widerspricht.}
\end{minipage}
\end{center}
\end{aufgabe}

\begin{aufgabe}
\boxtitel{Aufgabe 2.3 --- Wie viele Messungen reichen?}
Nach Satz 2.2: Wie viele Messungen braucht man, damit der Optimismus eines Polynoms vierten Grades höchstens $0{,}02$ beträgt? ($\sigma = 0{,}3$.)

\loesung
\schritt{Schritt 1 --- Die Parameterzahl}
Ein Polynom vierten Grades hat $p = 5$ Koeffizienten.

\schritt{Schritt 2 --- Die Bedingung}
\[ \frac{2\,\sigma^{2}\,p}{n} \ \le\ 0{,}02 \]
\[ \frac{2\cdot0{,}09\cdot5}{n} = \frac{0{,}9}{n} \ \le\ 0{,}02 \]
\[ n \ \ge\ \frac{0{,}9}{0{,}02} = 45 \]

\begin{center}
\begin{minipage}{0.9\textwidth}
\textbf{Fünfundvierzig Messungen --- also neun je Parameter.}
\end{minipage}
\end{center}

\schritt{Schritt 3 --- Die Faustregel}
\[ \frac{n}{p} \ \ge\ \frac{2\,\sigma^{2}}{\text{zulässiger Optimismus}} \]

\begin{center}
{\small
\renewcommand{\arraystretch}{1.48}
\begin{tabular}{@{}P{3.0cm}P{2.4cm}@{}}
\toprule
\textbf{zulässig} & $n/p$\\
\midrule
$0{,}05$ & $3{,}6$\\
$0{,}02$ & $9$\\
$0{,}01$ & $18$\\
\bottomrule
\end{tabular}}
\end{center}

\begin{center}
\begin{minipage}{0.88\textwidth}
\kursiv{In der klinischen Forschung gilt die Faustregel \enquote{zehn Ereignisse je Merkmal} für Vorhersagemodelle. \textbf{Hier sieht man, woher sie kommt} --- und dass sie nicht absolut ist, sondern vom Verhältnis von Rauschen zu zulässigem Fehler abhängt. Bei verrauschten Daten reichen zehn nicht; bei sehr genauen genügen weniger.}
\end{minipage}
\end{center}

\schritt{Schritt 4 --- Warum \enquote{Ereignisse} und nicht \enquote{Patienten}}
\begin{center}
\begin{minipage}{0.88\textwidth}
\kursiv{Bei einer seltenen Zielgröße --- zehn Prozent Sepsis --- trägt nur die kleinere Gruppe Information über den Unterschied. Tausend Patienten mit hundert Sepsisfällen sind für ein Modell mit zwanzig Merkmalen knapp: fünf Ereignisse je Merkmal. \textbf{Die wirksame Stichprobengröße ist die der seltenen Klasse}, und das wird in Studien regelmäßig übersehen.}
\end{minipage}
\end{center}
\end{aufgabe}

\begin{geschichte}
\boxtitelg{Ein Paradoxon, das niemand glauben wollte}
\Person{Charles Stein} (1920--2016) zeigte 1956, dass der Mittelwert nicht der beste Schätzer ist --- sobald man mindestens drei Größen gleichzeitig schätzt.

\begin{center}
\begin{minipage}{0.9\textwidth}
\kursiv{Schätzt man drei oder mehr unabhängige Mittelwerte, so gibt es einen Schätzer, der \kursiv{für jede} Wahl der wahren Werte einen kleineren Gesamtfehler hat als die einzelnen Mittelwerte: Man schrumpft alle zum gemeinsamen Durchschnitt hin. \textbf{Das gilt selbst dann, wenn die Größen nichts miteinander zu tun haben} --- etwa der Weizenpreis, die Zahl der Sepsisfälle und die Durchschnittsgröße von Schülern.}
\end{minipage}
\end{center}

\schritt{Die Aufnahme}
\begin{center}
\begin{minipage}{0.9\textwidth}
\kursiv{Das Ergebnis wurde zunächst für einen Fehler gehalten. \Person{Willard James} und Stein gaben 1961 den expliziten Schätzer an, der heute ihren Namen trägt. \textbf{Aufgabe 2.2 ist der eindimensionale Schatten davon} --- nur dass dort das Schrumpfen ein unbekanntes $\mu$ braucht, während James und Stein es aus den Daten selbst bestimmen, was erst ab drei Dimensionen gelingt.}
\end{minipage}
\end{center}

\schritt{Der Name}
Die Begriffe \enquote{Verzerrung} und \enquote{Varianz} sind alt; ihre Gegenüberstellung als Zielkonflikt des Lernens stammt von \Person{Stuart Geman}, \Person{Elie Bienenstock} und \Person{René Doursat} aus dem Jahr 1992.

\begin{center}
\begin{minipage}{0.9\textwidth}
\kursiv{Ihr Aufsatz handelte von neuronalen Netzen und war als Warnung gemeint: Große Netze hätten zu viel Varianz, um aus begrenzten Daten zu lernen. \textbf{Dreißig Jahre später ist das Bild verwickelter} --- Kapitel~14 zeigt, dass sehr große Netze die Kurve dieses Kapitels auf eine Weise verlassen, die niemand vorhergesagt hatte.}
\end{minipage}
\end{center}
\end{geschichte}

\begin{forschung}
\boxtitelp{Verzerrung und Varianz heute}
\begin{center}
{\small
\renewcommand{\arraystretch}{1.45}
\begin{tabular}{@{}P{2.8cm}P{5.4cm}@{}}
\toprule
\textbf{Doppelter Abstieg} & Jenseits der Interpolationsschwelle fällt der Testfehler wieder.\\
\textbf{Klinische Modelle} & Mindestgrößen für Vorhersagemodelle, die über die Zehnerregel hinausgehen.\\
\textbf{Ungleiche Streuung} & Modelle, die die Unsicherheit mitlernen.\\
\bottomrule
\end{tabular}}
\end{center}

\schritt{Zur ersten Zeile}
\begin{center}
\begin{minipage}{0.88\textwidth}
\kursiv{Die Tabelle dieses Kapitels endet bei Grad 9 --- genau dort, wo das Polynom die Daten interpoliert. \textbf{Was geschieht bei Grad 20 oder 100?} Die klassische Antwort lautet: noch schlimmer. Seit 2019 weiß man, dass unter bestimmten Umständen das Gegenteil eintritt: Der Testfehler erreicht an der Interpolationsschwelle ein Maximum und fällt danach wieder. Kapitel~14 behandelt das ausführlich.}
\end{minipage}
\end{center}

\schritt{Zur zweiten Zeile}
\begin{center}
\begin{minipage}{0.88\textwidth}
\kursiv{\Person{Richard Riley} und Mitarbeiter gaben ab 2019 Formeln für die Mindestgröße klinischer Vorhersagestudien an, die den Optimismus aus Satz 2.2 ausdrücklich begrenzen. \textbf{Sie ergeben oft deutlich mehr als zehn Ereignisse je Merkmal} --- und viele veröffentlichte Modelle erfüllen sie nicht.}
\end{minipage}
\end{center}
\end{forschung}

\begin{merksatz}
\boxtitel{Was Sie jetzt können}
\begin{center}
\renewcommand{\arraystretch}{1.5}
\begin{tabular}{@{}P{3.0cm}P{5.6cm}@{}}
\toprule
\textbf{Satz 2.1} & Fehler $=$ Rauschen $+$ Verzerrung$^{2}$ $+$ Varianz\\
\textbf{Unteranpassung} & zu starr, Verzerrung überwiegt\\
\textbf{Überanpassung} & zu frei, Varianz überwiegt\\
\textbf{Satz 2.2} & Optimismus $2\sigma^{2}p/n$\\
\textbf{Mehr Daten} & senken die Varianz, nicht die Verzerrung\\
\bottomrule
\end{tabular}
\end{center}

\begin{center}
\begin{minipage}{0.88\textwidth}
\kursiv{Im Leitfall kannten wir die wahre Kurve und konnten den Testfehler ausrechnen. Das nächste Kapitel fragt, wie man ihn schätzt, wenn man sie nicht kennt --- und wie man dabei betrügt, ohne es zu merken.}
\end{minipage}
\end{center}
\end{merksatz}
